\chapter{Technical Background}\label{chap:theory}

This chapter covers the theoretical foundations of the key methods I built on.
My aim is not to reproduce full mathematical proofs but to explain clearly what each
method does, present the central equations that govern its behaviour, and make the
connection to how I applied it in this project explicit.

The chapter is organised around five areas:
(i)~the transformer architecture and self-attention mechanism that FinBERT is built on;
(ii)~BERT's pre-training procedure and FinBERT's financial adaptation;
(iii)~LIME's local surrogate model for explainability;
(iv)~SHAP and its game-theoretic grounding; and
(v)~the statistical correlation and causal inference methods used in the correlation engine.

% ----------------------------------------------------------
\section{The Transformer Architecture}
% ----------------------------------------------------------

The transformer \citep{vaswani2017attention} replaced recurrent and convolutional sequence
models with a purely attention-based architecture.
The key insight is that any two positions in a sequence can attend directly to one another,
removing the sequential bottleneck that made RNNs slow to train and prone to forgetting
over long contexts.  Figure~\ref{fig:transformer_encoder} shows the structure of a single encoder layer.

\begin{figure}[H]
\centering
\begin{tikzpicture}[
  node distance=0.75cm,
  every node/.style={font=\small},
  mainbox/.style={draw, rounded corners=4pt, minimum width=7cm,
                  minimum height=0.80cm, fill=#1, align=center,
                  inner sep=7pt, text width=6.6cm},
  normbox/.style={draw, rounded corners=4pt, minimum width=7cm,
                  minimum height=0.65cm, fill=gray!12, align=center,
                  inner sep=5pt, font=\small},
  arr/.style ={->, >=Stealth, thick},
  skip/.style={->, >=Stealth, thick, dashed, blue!55}]

\node[mainbox=teal!18]   (embed) {Token Embeddings + Positional Encoding};
\node[mainbox=blue!15, below=of embed]  (mha)  {Multi-Head Self-Attention\\[-2pt]\scriptsize $h{=}12$ heads, $d_k{=}64$ per head};
\node[normbox,  below=of mha]           (ln1)  {Add \& Layer Norm};
\node[mainbox=purple!13, below=of ln1]  (ffn)  {Position-wise Feed-Forward Network\\[-2pt]\scriptsize $\mathrm{FFN}(\mathbf{x})=\mathrm{ReLU}(\mathbf{x}W_1+b_1)W_2+b_2$};
\node[normbox,  below=of ffn]           (ln2)  {Add \& Layer Norm};
\node[mainbox=teal!18, below=of ln2]    (outp) {Encoder Output\\[-2pt]\scriptsize ($\times\,12$ stacked layers in BERT-base)};

\draw[arr] (embed) -- (mha);
\draw[arr] (mha)   -- (ln1);
\draw[arr] (ln1)   -- (ffn);
\draw[arr] (ffn)   -- (ln2);
\draw[arr] (ln2)   -- (outp);

%% Residual connections (left rail, away from the main vertical flow)
\draw[skip] (embed.west) -- ++(-1.15,0)
    node[left, font=\scriptsize\itshape, text=gray!80] {residual}
    |- (ln1.west);
\draw[skip] (ln1.west) -- ++(-1.15,0)
    node[left, font=\scriptsize\itshape, text=gray!80] {residual}
    |- (ln2.west);
\end{tikzpicture}
\caption{Structure of a single transformer encoder layer as used in BERT and FinBERT.  Both the multi-head attention and feed-forward sub-layers are wrapped in a residual (skip) connection followed by layer normalisation, preserving gradient flow through the 12-layer stack.}
\label{fig:transformer_encoder}
\end{figure}

% ............................................................
\subsection{Positional Encoding}

Because the attention mechanism has no inherent sense of word order, position information
must be injected explicitly.
Vaswani et al.\ achieve this with a fixed encoding that uses sinusoidal functions of token
position $t$ and embedding dimension index $i$:
\begin{align}
  PE_{(t,\,2i)}   &= \sin\!\left(\frac{t}{10000^{2i/d_{\text{model}}}}\right) \\[4pt]
  PE_{(t,\,2i+1)} &= \cos\!\left(\frac{t}{10000^{2i/d_{\text{model}}}}\right)
\end{align}
where $d_{\text{model}} = 768$ in BERT-base.
The sinusoidal choice means that each position gets a unique, smoothly varying signature
across the embedding dimensions, and the model can learn to attend to relative positions
rather than absolute ones.
The positional encoding is added element-wise to the token embedding before the first
encoder layer.

% ............................................................
\subsection{Scaled Dot-Product Attention}

The central operation of the transformer is \emph{scaled dot-product attention}.
Inputs are projected into three matrices: queries $\mathbf{Q}$, keys $\mathbf{K}$,
and values $\mathbf{V}$. The output is a weighted sum of the values, where the
weights reflect how relevant each key is to each query:
\begin{equation}
  \text{Attention}(\mathbf{Q},\mathbf{K},\mathbf{V})
  = \text{softmax}\!\left(\frac{\mathbf{Q}\mathbf{K}^{\!\top}}{\sqrt{d_k}}\right)\mathbf{V}.
  \label{eq:attention}
\end{equation}

Each entry of $\mathbf{Q}\mathbf{K}^{\!\top}$ scores how strongly one position
attends to another.
After the softmax, each row becomes a probability distribution over all positions, and
multiplying by $\mathbf{V}$ produces a weighted average of value vectors.
The $\sqrt{d_k}$ scaling prevents the dot products from growing so large that the
softmax saturates into near-zero gradients, a practical issue that appeared
consistently in early training runs before scaling was introduced.

% ............................................................
\subsection{Multi-Head Attention}

Rather than applying one attention function to the full representation,
the transformer runs $h$ parallel \emph{attention heads}, each projecting into a separate
lower-dimensional subspace:
\begin{equation}
  \text{MultiHead}(\mathbf{Q},\mathbf{K},\mathbf{V})
  = \text{Concat}(\text{head}_1,\ldots,\text{head}_h)\,\mathbf{W}^O
  \label{eq:mha}
\end{equation}
where each $\text{head}_i = \text{Attention}(\mathbf{Q}\mathbf{W}_i^Q,\,
\mathbf{K}\mathbf{W}_i^K,\, \mathbf{V}\mathbf{W}_i^V)$ and the results are concatenated
and projected back via $\mathbf{W}^O$.
BERT-base uses $h = 12$ heads.

The motivation is that a single attention distribution must simultaneously resolve
multiple linguistic tasks.
Multiple heads allow the model to specialise: some learn local phrase-level dependencies,
others capture long-range entity references.
In financial text this matters because the sentiment-bearing term (e.g.\ ``surged'')
may appear several tokens away from the company name it modifies.

% ............................................................
\subsection{Feed-Forward Sub-layers and Residual Connections}

After the attention sub-layer, each encoder layer applies a two-layer feed-forward network
independently at every position:
\begin{equation}
  \text{FFN}(\mathbf{x}) = \max(0,\;\mathbf{x}\mathbf{W}_1 + \mathbf{b}_1)\,\mathbf{W}_2 + \mathbf{b}_2.
  \label{eq:ffn}
\end{equation}
Both the attention and FFN sub-layers are wrapped in a residual connection followed by
layer normalisation, which stabilises training across the 12-layer stack by preserving
a gradient path directly to earlier layers.

% ----------------------------------------------------------
\section{BERT and FinBERT}
% ----------------------------------------------------------

% ............................................................
\subsection{BERT Pre-training}

BERT \citep{devlin2019bert} pre-trains on two self-supervised objectives over large
unlabelled text corpora.

\textbf{Masked Language Modelling (MLM)} randomly selects 15\% of input tokens,
replacing 80\% of those with a \texttt{[MASK]} token, 10\% with a random vocabulary
token, and leaving 10\% unchanged.
The model is trained to predict the original token at each masked position from the
surrounding context.
The 10\%/10\% split is deliberate: if every selected token became \texttt{[MASK]},
the model would only learn to handle masked inputs, creating a mismatch at inference
time where \texttt{[MASK]} never appears.

\textbf{Next Sentence Prediction (NSP)} presents the model with pairs of sentences
and asks it to predict whether the second genuinely follows the first.
A special \texttt{[CLS]} token prepended to each input is used for the binary
classification.
RoBERTa \citep{liu2019roberta} later showed NSP to be unnecessary, but it is part of
the BERT setup that FinBERT inherits.

% ............................................................
\subsection{FinBERT Fine-tuning for Sentiment Classification}

Araci \citep{araci2019finbert} takes the pre-trained BERT-base weights, continues
pre-training on the Reuters TRC2 corpus and the Financial PhraseBank
\citep{malo2014financialphrasebank}, and then fine-tunes for three-class sentiment
classification (positive, neutral, negative).

Fine-tuning adds a single linear layer on top of the \texttt{[CLS]} token's final
representation $\mathbf{h}_{\texttt{[CLS]}}$:
\begin{equation}
  \hat{\mathbf{y}}
  = \text{softmax}\!\left(\mathbf{W}_c\,\mathbf{h}_{\texttt{[CLS]}} + \mathbf{b}_c\right),
  \label{eq:finbert_clf}
\end{equation}
where $\mathbf{W}_c \in \mathbb{R}^{3 \times 768}$.
Training minimises the standard cross-entropy loss between the predicted distribution
$\hat{\mathbf{y}}$ and the one-hot ground-truth label.

Domain pre-training matters because financial vocabulary is highly context-dependent.
The phrase ``the company cut its losses'' carries a positive connotation in a trading
context (implying a controlled exit), whereas a general-purpose model would classify
``cut'' and ``losses'' together as negative.
I chose FinBERT specifically because I needed financial understanding baked into the
representations before fine-tuning, not retrofitted afterwards.

% ----------------------------------------------------------
\section{Local Interpretable Model-Agnostic Explanations (LIME)}
% ----------------------------------------------------------

% ............................................................
\subsection{Core Idea}

LIME \citep{ribeiro2016lime} explains the prediction of any black-box model $f$ at a
specific input $x^*$ by fitting a simple, interpretable surrogate model in the local
neighbourhood of $x^*$.
Its \emph{model-agnostic} property means the same method applies regardless of whether
$f$ is FinBERT, a decision tree, or any other classifier
\citep{doshivelez2017interpretable}.  The end-to-end workflow is shown in Figure~\ref{fig:lime_process}.

\begin{figure}[H]
\centering
\resizebox{\textwidth}{!}{%
\begin{tikzpicture}[
  node distance=1.2cm and 1.6cm,
  every node/.style={font=\small},
  lmbox/.style={draw, rounded corners=5pt, fill=#1, minimum width=3.6cm,
                minimum height=1.5cm, align=center, inner sep=9pt,
                text width=3.2cm},
  num/.style={circle, draw, fill=white, minimum size=0.6cm,
              font=\small\bfseries, inner sep=0pt},
  arr/.style={->, >=Stealth, very thick, line width=1.2pt}]

%% Top row: steps 1–3
\node[lmbox=teal!18]   (input)   {\textbf{1.\ Input text} $x^*$\\[4pt]
  \scriptsize Raw financial sentence fed to the explainer};
\node[lmbox=blue!14,  right=of input]   (perturb) {\textbf{2.\ Perturbation}\\[4pt]
  \scriptsize $\sim$200 randomly masked variants of $x^*$ generated};
\node[lmbox=orange!16, right=of perturb] (query)   {\textbf{3.\ FinBERT} $f(z)$\\[4pt]
  \scriptsize Each variant scored; label + probability recorded};

%% Bottom row: steps 4–5 (right to left to form a U-shape)
\node[lmbox=purple!14, below=of query]   (fit)     {\textbf{4.\ Surrogate fitting}\\[4pt]
  \scriptsize Proximity-weighted linear regression on perturbed outputs};
\node[lmbox=green!16,  left=of fit]      (output)  {\textbf{5.\ Token weights} $w_i$\\[4pt]
  \scriptsize Coefficient of each token = its influence on the prediction};
\node[lmbox=teal!10,   left=of output]   (heatmap) {\textbf{Output: heatmap}\\[4pt]
  \scriptsize Rendered in the dashboard's explainability panel};

%% Step number badges – centred exactly on each box's north-west corner so the
%% disc straddles the corner (half inside / half outside) and never overlaps text.
\foreach \n/\nd in {1/input,2/perturb,3/query,4/fit,5/output}{
  \node[num] at (\nd.north west) {\n};
}

%% Arrows
\draw[arr] (input)   -- (perturb);
\draw[arr] (perturb) -- (query);
\draw[arr] (query)   -- (fit);
\draw[arr] (fit)     -- (output);
\draw[arr] (output)  -- (heatmap);
\end{tikzpicture}}
\caption{LIME explanation process applied to financial text classification.  The explainer generates perturbed variants of the input, queries the black-box FinBERT model on each, and fits a proximity-weighted linear surrogate to recover per-token importance weights.  These weights are displayed as a colour-coded heatmap in the dashboard.}
\label{fig:lime_process}
\end{figure}

% ............................................................
\subsection{The LIME Objective}

LIME searches for an interpretable model $g$ (here, a linear model) that balances
fidelity to $f$ locally against simplicity:
\begin{equation}
  \xi(x^*)
  = \arg\min_{g \;\in\; \mathcal{G}}
    \;\mathcal{L}(f,\,g,\,\pi_{x^*}) + \Omega(g)
  \label{eq:lime_obj}
\end{equation}
where $\mathcal{L}(f, g, \pi_{x^*})$ is the weighted prediction error between $f$ and $g$
over perturbed samples near $x^*$, $\Omega(g)$ penalises the complexity of $g$, and
$\pi_{x^*}(z) = \exp(-D(x^*, z)^2 / \sigma^2)$ is a proximity kernel that gives more
weight to samples close to $x^*$.

% ............................................................
\subsection{Perturbation and Surrogate Fitting for Text}

For a text input, LIME treats each word as a binary feature (present or absent).
It generates a set of perturbed versions of $x^*$ by randomly masking tokens, queries $f$
on each, and uses the resulting predictions as labels to fit the weighted linear surrogate.
The coefficients of the fitted model are token importance scores: a positive weight on a
word means that word pushed the prediction toward the target class; a negative weight means
it pushed against it.
This produces the token-level heatmap shown in the dashboard's explainability panel.

The key limitation is that a linear surrogate cannot capture non-local interactions,
for example negation that spans several tokens, which BERT's attention mechanism
handles naturally.
LIME's explanation is therefore an approximation of $f$'s behaviour near $x^*$, not a
global account of the model.
I accepted this trade-off because LIME's cost scales with the number of perturbations
sampled (around 200 for the text lengths I was working with), making it fast enough to
serve from a live API endpoint.

% ----------------------------------------------------------
\section{SHAP: SHapley Additive Explanations}
% ----------------------------------------------------------

% ............................................................
\subsection{Game-theoretic Foundation and the Shapley Value}

SHAP \citep{lundberg2017shap} grounds feature attribution in cooperative game theory.
Each feature $x_i$ is treated as a ``player'', and the model's output is the
``prize'' to be divided among the players.
The \emph{Shapley value} $\phi_i$ measures the average marginal contribution of feature
$i$ across all possible orderings in which features can be added to a coalition:
\begin{equation}
  \phi_i
  = \sum_{S \;\subseteq\; F \setminus \{i\}}
    \frac{|S|!\;(|F| - |S| - 1)!}{|F|!}
    \;\bigl[f(S \cup \{i\}) - f(S)\bigr]
  \label{eq:shapley}
\end{equation}
where $F$ is the full set of features, $S$ is a coalition not containing $i$,
$f(S)$ is the model's expected output given only the features in $S$, and the
combinatorial factor weights each coalition by how many orderings it corresponds to.

The Shapley value is the unique attribution that satisfies a set of fairness axioms,
including that attributions sum exactly to the model's output and that features with
equal contribution receive equal attribution.

% ............................................................
\subsection{Computational Cost and the Choice of LIME at Runtime}

Exactly evaluating~\eqref{eq:shapley} requires considering all $2^{|F|}$ subsets.
For a 128-token input this is entirely intractable.
KernelSHAP \citep{lundberg2017shap} approximates Shapley values by solving a
specially weighted regression, conceptually similar to LIME but with a different
kernel that is proven to recover the Shapley values in expectation.
Even with this approximation, a single SHAP explanation requires several hundred model
evaluations, versus LIME's $\sim$200 perturbations on comparable inputs.
This is the practical reason I served LIME from the real-time API and reserved SHAP
for offline notebook analysis.

% ----------------------------------------------------------
\section{Correlation and Causal Inference Methods}
% ----------------------------------------------------------

% ............................................................
\subsection{Pearson and Spearman Correlation}

Given two time series of length $T$, the Pearson correlation coefficient measures linear
co-movement:
\begin{equation}
  r = \frac{\sum_{t=1}^T (x_t - \bar{x})(y_t - \bar{y})}
           {\sqrt{\sum_{t=1}^T (x_t - \bar{x})^2}\;\;\sqrt{\sum_{t=1}^T (y_t - \bar{y})^2}},
  \label{eq:pearson}
\end{equation}
$r$ lies in $[-1, 1]$, with $+1$ indicating perfect positive linear co-movement.
Its limitation is that it assumes a bivariate normal distribution, which financial
returns do not satisfy (heavy tails, occasional extreme values).

Spearman's $r_s$ addresses this by applying the same formula to the \emph{ranks} of the
observations rather than the raw values, which removes distributional assumptions.
When there are no tied ranks it simplifies to:
\begin{equation}
  r_s = 1 - \frac{6\sum_{t=1}^T d_t^2}{T(T^2 - 1)},
  \label{eq:spearman}
\end{equation}
where $d_t$ is the difference in ranks of $x_t$ and $y_t$.
I report both measures so that the distance between them indicates how much of any
apparent relationship is an artefact of outliers or non-normality.

% ............................................................
\subsection{Rolling Window Correlation}

A single static correlation coefficient masks how the relationship changes over time;
for instance, sentiment--price coupling may strengthen during high-volatility periods and
weaken in calmer markets.
Rolling correlation applies the Pearson formula within a sliding window of width $W$,
producing a time series of coefficients $\{r_t^{(W)}\}$ that reveals these temporal
shifts.
I implemented this using \texttt{pandas.DataFrame.rolling().corr()} with a 30-day window,
which balanced responsiveness to regime changes against the noise introduced by short
windows.

% ............................................................
\subsection{Granger Causality}

Granger causality \citep{granger1969} asks: do past values of series $X$ improve
predictions of $Y$ beyond what $Y$'s own history alone provides?
It is important to note that this is \emph{predictive precedence}, not physical causation
in any deeper sense.

\paragraph{Stationarity.}
Both series must be stationary before the test is valid; non-stationary series can
generate spurious correlations \citep{engle1987cointegration}.
Sentiment scores are bounded and approximately stationary; stock prices are converted to
log-differenced returns $r_t = \log(P_t / P_{t-1})$ to remove the price trend.

\paragraph{Model comparison.}
For lag order $p$ (chosen by the Akaike Information Criterion), two autoregressive
models are fitted.
The \emph{unrestricted} model regresses $Y$ on its own past values and the past values
of $X$:
\begin{equation}
  y_t
  = \alpha_0
    + \sum_{k=1}^{p} \alpha_k\, y_{t-k}
    + \sum_{k=1}^{p} \beta_k\, x_{t-k}
    + \varepsilon_t.
  \label{eq:var_unrest}
\end{equation}
The \emph{restricted} model uses only $Y$'s own history:
\begin{equation}
  y_t
  = \alpha_0
    + \sum_{k=1}^{p} \alpha_k\, y_{t-k}
    + \varepsilon_t^R.
  \label{eq:var_rest}
\end{equation}

\paragraph{The F-test.}
The null hypothesis $H_0\colon \beta_1 = \cdots = \beta_p = 0$ is tested with:
\begin{equation}
  F = \frac{(RSS_R - RSS_U)\,/\,p}{RSS_U\,/\,(T - 2p - 1)}
  \;\sim\; F\!\left(p,\;T - 2p - 1\right),
  \label{eq:granger_f}
\end{equation}
where $RSS_R$ and $RSS_U$ are the residual sums of squares of the restricted and
unrestricted models.
A $p$-value below 0.05 means the added sentiment lags significantly reduce the
prediction error for returns, and $X$ is said to Granger-cause $Y$.

I found the results varied considerably by ticker.
Sentiment appeared to Granger-cause returns for some large-cap, news-driven equities
but showed no significant effect for others.
This is consistent with theory: more heavily traded, high-coverage stocks are more
sensitive to retail commentary than smaller or less liquid names where sentiment volume
is too low for the test to have power.
The inconsistency is itself a finding about which assets are most affected by public
sentiment.

% ----------------------------------------------------------
\section{Summary}
% ----------------------------------------------------------

This chapter has covered the theoretical basis for five method families used in the project.
The transformer's attention mechanism and FinBERT's domain pre-training explain why the
model performs well on financial text; LIME and SHAP each provide a principled route to
explaining individual predictions, with different cost and fidelity trade-offs; and
Pearson, Spearman, rolling correlation, and the Granger F-test together give a layered
view of the sentiment--price relationship that no single measure could provide alone.

Understanding these foundations directly shaped how I built and debugged the system.
Knowing that LIME approximates $f$ locally with a linear model, for instance, helped me
interpret cases where the token heatmap contradicted my intuition, and the disagreement was
not a bug in my code but a consequence of the linear surrogate breaking down in a
non-linear region of the decision boundary.
The next chapter describes how these methods were combined into a working system.
